\section{Task 7 experimental results}
The supplied dataset contains 7112 character images. The development portion has 5334 images and the held-out test has 1778 images. Exact duplicate groups remain within one partition.
\subsection{Held-out classification}
\begin{center}\begin{tabular}{lrr}\toprule
Method & Accuracy & Macro F1 \\\midrule
CNN & 99.38\% & 0.99382 \\
SVM & 99.55\% & 0.99550 \\
\bottomrule\end{tabular}\end{center}
\subsection{Chip-image inference}
The reference string is \texttt{HD44780A00}.
\begin{center}\begin{tabular}{lllr}\toprule
Connectivity & Method & Prediction & Position accuracy \\\midrule
4 & CNN & \texttt{HD44780A00} & 100.00\% \\
4 & SVM & \texttt{HD44780A00} & 100.00\% \\
8 & CNN & \texttt{HD44780A00} & 100.00\% \\
8 & SVM & \texttt{HD44780A00} & 100.00\% \\
\bottomrule\end{tabular}\end{center}
Model configurations and epoch budgets were chosen using internal validation, followed by a fixed-budget refit on all development images. These measurements use one grouped split and one seed. The SVM uses hand-written HOG and a linear squared-hinge objective optimized with Adam. Both optimizers use PyTorch. The chip image supplies feedback during SVM refinement, and its transcription score describes this development example.
The experiment tables are generated from the saved JSON records. The group should develop its report discussion from its own assessment of these results.
